\section{The unitary Weil normal-core gate}
\label{sec:unitary-weil-core}

\paragraph{Status and scope.}
This is a new author derivation from published Weil-character and
automorphism inputs, with complete retained finite endpoint tables.
The author has reread the entire proof and executed both the complete
GAP scan and a separate exact Python reconstruction; all pass.
An independent whole-proof, primary-source and complete arithmetic
audit by /root/imprimitive has passed. It concerns
ordinary unitary Weil representations, not arbitrary irreducible
representations of unitary groups. The transitivity hypothesis below
is explicit; it is the one supplied by D-BI-19 in the main primitive,
tensor-indecomposable gate. No equality of the original and inducing
character fields is assumed.

Write $\nu(T)=a-\max_\lambda\dim\ker(T-\lambda)$ on a space of
dimension $a$. A scalar has $\nu=0$; a genuine nonidentity scalar
inertia has ordinary moving rank $a$. These are different quantities.

\begin{lemma}[D-BI-24a: unitary Weil characters and stabilizers]
\label{D-BI-24a}
Let $n\ge3$, let $q=p^f$, and suppose $(n,q)\ne(3,2)$.
Put $\epsilon=(-1)^n$. The ordinary unitary Weil constituents
$W_i$, $0\le i\le q$, restrict to pairwise distinct irreducible
representations of $L=\operatorname{SU}_n(q)$. Their degrees are
\[
 a_0=\frac{q^n+q\epsilon}{q+1},\qquad
 a_i=\frac{q^n-\epsilon}{q+1}\quad(i\ne0).                  \tag{UW.1}
\]
Every nonscalar $g\in\operatorname{GU}_n(q)$ satisfies
\[
 |\chi_i(g)|\le B:=\frac{q^{n-1}+2q-1}{q+1},\qquad
 a_i\ge A:=\frac{q^n-q}{q+1}.                              \tag{UW.2}
\]
Complex conjugation sends $\chi_i|_L$ to $\chi_{-i}|_L$,
and entrywise $p$-Frobenius sends its label to $pi$ modulo $q+1$.

For $q=2,3,5$, the full linear normalizer of a nonselfdual
constituent is contained in the scalar multiples of its
$\operatorname{GU}_n(q)$-extension. The same assertion holds
for every $i\ne0$ when $(n,q)=(3,4)$.
For each such nonscalar normalizer $T$,
\[
                      \nu(T)\ge(a_i-B)/2.                  \tag{UW.3}
\]
\end{lemma}

\begin{proof}
The published input is G\'erardin's total unitary Weil representation
$\Omega$ and its exact character formula
\[
 \operatorname{Tr}\Omega(g)
       =(-1)^n(-q)^{\dim_{\mathbf F_{q^2}}\ker(g-1)};
                                                               \tag{UW.4}
\]
see \cite[Theorem4.9.2 and its complete proof, pp.~99--100]{BI:GerardinLinear}.
This input includes even characteristic: unitary \S3.1 on p.~70
starts with an arbitrary finite field, and Theorem3.3(a\textnormal{'''})
on p.~73 explicitly constructs the even-$q$ representation.
The odd-characteristic restriction in that article belongs to the
symplectic \S2 and Theorem4.9.1, not to Theorem4.9.2.
The central-character projector formula is also explicitly recorded
in \cite[\S4.3, (4.3)--(4.4), author-version p.~16]{BI:GLOST2018};
the ordinary degrees and extensions are stated in
\cite[p.~293]{BI:GMST2002}.

Here is a direct check of the restriction assertion used below.
The scalar center of $\operatorname{GU}_n(q)$ is cyclic of order
$q+1$. Projecting $\Omega$ to its characters gives $q+1$ nonzero
spaces, with degrees (UW.1), by evaluating (UW.4) at scalar matrices.
Furthermore, $|\operatorname{Tr}\Omega(g)|^2=q^{2\dim\ker(g-1)}$
is the permutation character of $L$ on its natural vector space.
This action has exactly $q+1$ orbits: zero, nonzero isotropic
vectors, and one orbit for each of the $q-1$ nonzero norm values.
Witt extension gives transitivity for the full unitary group.
To adjust the determinant to one while fixing a target vector,
use a nondegenerate complement of dimension $n-1$ in the
anisotropic case, or a nondegenerate complement of dimension
$n-2$ to a hyperbolic plane in the isotropic case. Either
complement contains an anisotropic line on which any scalar
of norm one may act. Thus the determinant adjustment is available
for $n\ge3$, proving the same orbit assertion for $L$.
Consequently $\langle\Omega|_L,\Omega|_L\rangle=q+1$.
There are already $q+1$ nonzero central summands; positivity of
character inner products forces each restriction to be irreducible
and distinct from all the others.

Fix a primitive $(q+1)$st root $\rho$ in $\mathbf F_{q^2}$ and
a corresponding complex root $\zeta$. With a choice of label sign,
the projector formula is
\[
 \chi_i(g)=\frac{(-1)^n}{q+1}
       \sum_{j=0}^{q}\zeta^{ij}(-q)^{d_j(g)},\qquad
 d_j(g)=\dim\ker(g-\rho^j).                                \tag{UW.5}
\]
For a nonscalar natural matrix, $0\le d_j\le n-1$ and
$\sum_jd_j\le n$. The convex sum $\sum_jq^{d_j}$ is at
most its value at $(n-1,1,0,\ldots,0)$: move one unit from
a smaller positive entry to a larger entry until an entry
reaches $n-1$, and then concentrate the remaining unit.
That value is $q^{n-1}+q+(q-1)$, proving (UW.2).
Formula (UW.5) also proves the two label rules directly.

For precision, the published automorphism input is
\cite[\S2.2, Lemma6, p.~114]{BI:YangWuRevinVdovin2023}:
\[
 \operatorname{Aut}(\operatorname{PSU}_n(q))
      =\operatorname{PGU}_n(q)\rtimes\langle\sigma\rangle,
 \qquad |\sigma|=2f,                                       \tag{UW.6}
\]
where $\sigma$ is induced by entrywise $p$-Frobenius.
Only the standard simple unitary groups specified here are used.
For $n\ge3$ the nonsimple pair $(3,2)$ has already been excluded.
In particular the small groups $\operatorname{PSU}_3(3)$,
$\operatorname{PSU}_4(2)$, $\operatorname{PSU}_5(2)$,
$\operatorname{PSU}_6(2)$ and $\operatorname{PSU}_3(4)$,
$\operatorname{PSU}_3(5)$ all lie within this stated domain.
If a linear normalizer induces an automorphism on the represented
central quotient of $L$, take a standard lift $\alpha$ from
(UW.6) to $L$. Its action and conjugation by the normalizer agree
projectively. Their ratio is a scalar character of $L$, hence
trivial because $L$ is perfect. This also shows that the lift
preserves the kernel of the represented constituent whenever
that normalizer exists. No unproved characteristic-kernel claim
about an exceptional Schur multiplier is required.

All inner-diagonal automorphisms stabilize the constituent,
since it extends to $\operatorname{GU}_n(q)$. For $q=2,3,5$
the nontrivial field quotient sends $i$ to $-i$; it cannot
stabilize a nonselfdual constituent, by the distinctness just
proved. When $q=4$, multiplication by $2$ has order four on
every nonzero label modulo five, so no nonidentity field
quotient stabilizes such a label. Schur's lemma now makes
every allowed normalizer a scalar multiple of the known
unitary-group operator. A nonscalar represented operator
has a nonscalar natural lift. Finally, if an eigenspace has
dimension $m$, rotating its eigenvalue to one gives
$|\operatorname{Tr}(T)|\ge2m-a_i$. This and (UW.2) give (UW.3).
\end{proof}

\begin{lemma}[D-BI-24b: the three exact small endpoints]
\label{D-BI-24b}
For the following ordinary constituents, every nonscalar
finite-order linear normalizer of the represented simple
core has the indicated support bound:
\[
\begin{array}{c|c|c}
\text{simple core}&\text{constituent degrees}&\nu(T)\text{ lower bound}\\\hline
\operatorname{PSU}_3(3)&6,7&a/4\\
\operatorname{PSU}_5(2)&10&a/5\\
\operatorname{PSU}_6(2)&22&a/4.
\end{array}                                                  \tag{UW.7}
\]
The degree-ten constituent is selfdual. The degree-twenty-two
constituent is $W_0$ and is trivial on the center of
$\operatorname{SU}_6(2)$.
\end{lemma}

\begin{proof}
These are complete finite character-table computations from
published ATLAS inputs, not inferred normalizer bounds from
samples. The retained seven tables, selected rows and exact
nonscalar-entry counts are
\[
\begin{array}{c|c|c|c}
\text{table}&\text{rows}&\text{entries}&\min\nu\\\hline
\texttt{U3(3)}&2,3,4,5&52&2,3,2,2\\
\texttt{U3(3).2}&3,4,5,6&60&2,2,3,3\\
\texttt{U5(2)}&2&46&2\\
\texttt{U5(2).2}&3,4&84&2,2\\
\texttt{U6(2)}&2&45&6\\
\texttt{U6(2).2}&3,4&128&6,6\\
\texttt{U6(2).3}&4,5,6&339&6,6,6.
\end{array}                                                  \tag{UW.8}
\]
All characters, class sizes and orders, prime power maps,
Frobenius--Schur indicators, source strings, and core fusions
are retained in \texttt{unitary\_weil\_inputs.json}.
The exporter \texttt{unitary\_weil/export\_and\_scan.g}
checks full row orthogonality. The separate exact Python
checker reselects rows by pointwise restriction equality,
checks the normal fusions and cyclic quotients, and reconstructs
every eigenspace multiplicity from all powers:
\[
 m_b=\frac1d\sum_{j=0}^{d-1}\chi(g^j)\zeta_d^{-bj}.
                                                               \tag{UW.9}
\]
It checks integrality, nonnegativity and total degree. Thus
all seventeen selected rows and all754 nonscalar spectra
are checked. The full tables retain their published status;
orthogonality is a consistency check and does not replace
the published identification of the groups and tables.
The retained input SHA-256 is
\texttt{0c84735f987b0bea3a9fea197002d4cb28af2cdae9a165bad8bce3bf6f91a117};
the Python program and actual output are
\texttt{check\_unitary\_weil.py} and
\texttt{check\_unitary\_weil.json}. The first GAP attempt
stopped on an unsupported metadata operation for an extension
table. That metadata call was restricted to the simple core;
no character data changed, and the complete rerun passed.

The simple outer groups in these cases are respectively
$C_2,C_2,S_3$. This follows either from the retained official
extension metadata or from (UW.6): for the third group the
diagonal group is $C_3$, inverted by the field involution.
The specified extensions cover every cyclic subgroup of
each outer group up to conjugacy. For an actual normalizer,
adjust by an inner operator to match a chosen generator
$t$ of one such fixed extension $E/S$. If its order modulo
$S$ is $d$, then $T^d\rho(t^d)^{-1}$ is scalar by Schur's
lemma. Rescale $T$ by a $d$th root of the inverse scalar.
The conjugation and power relations then give a representation
of the fixed $E$, irreducible because its restriction to $S$
is irreducible. It is one of the exactly selected rows.
Putting back the inner operator still gives a checked element
of $E$, and scalar rescaling does not change $\nu$.

For the degree-ten and degree-twenty-two cases, the unique
rows in (UW.8) are exactly $W_0$ by (UW.1); (UW.5) makes
their characters real. In degree22 the central projector
is the trivial one, so the representation factors through
the simple group as stated. This proves every assertion.
\end{proof}

\begin{theorem}[D-BI-24: reduction of all transitive unitary Weil cores]
\label{D-BI-24}
Consider an inducing image $J\subset\operatorname{GL}(V)$
of an actual rank-nine Weil sixfold finite-cover source.
Suppose its character field contains the original imaginary
quadratic field $K$. Suppose $J$ has an irreducibly acting
normal central product $N=N_1\cdots N_t$ with all simple
quotients isomorphic to $\operatorname{PSU}_n(q)$,
$n\ge3$, $(n,q)\ne(3,2)$, and with the factor representations
given by ordinary unitary Weil constituents. Suppose $J$
acts transitively on the factors.

Then necessarily
\[
                   (n,q)=(4,2),\qquad t=1,\qquad
                   \dim V\in\{5,6\}.                       \tag{UW.10}
\]
For the eight non-Eisenstein class-number-one fields of the
main nonsplit problem, with the actual original polarization
nonsplit, only degree five remains. It is subject to the
already proved rank-five inducing restrictions; no new
existence or dominance claim is made for it.
\end{theorem}

\begin{proof}
Write $a=\dim W_i$ for the common factor degree, so $e=\dim V=a^t$.
Transitivity makes the constituent labels automorphically
related and in particular makes selfduality uniform.
The intrinsic tensor matrix algebras of the simple factors
are permuted by $J$. A factor-preserving normalizer is a
tensor product of factor intertwiners. The maximum eigenspace
fraction of a tensor product is at most that of any nonscalar
factor. A nontrivial weighted factor cycle has the stronger
quarter support bound from D-SPOR-03. The cycle product returns
to the same constituent; this permits identifications between
different automorphic or dual labels without allowing an
incompatible field automorphism on a fixed constituent.

The actual intermediate-cover Euler identity is
$D=6+2e(1-h)$. Here $a\ge5$, so genus $h\ge2$ is impossible.
For $h=1$, at most six genuine branch points remain; the
marked elliptic-base dimension is at most six (and is one
in the unbranched case). The actual period variation factors
through that finite Hurwitz cover and cannot have rank nine.
Thus a rank-nine source necessarily has
\[
                h=0,\qquad r\ge12,\qquad D=2e+6.            \tag{UW.11}
\]
This is only a necessary consequence of the actual variation
factorization, not a claim of dominance from branch dimension.

For the generic trace estimate put $c=B/a$. On the central
product, every nonscalar element has character ratio at
most $c$. D-BI-18 therefore gives
\[
                       e(1-3B/a)\le6.                      \tag{UW.12}
\]
When all factors are selfdual, their tensor product has a
unique invariant nondegenerate bilinear form line, symmetric
or alternating. Every element of $J$ acts on it by a multiplier.
The multiplier is nontrivial: otherwise the finite character
of $J$ is real, contradicting $K\subseteq\mathbf Q(\chi_J)$.
The generating product-one inertia tuple has at least two
nontrivial multipliers; their fixed spaces are isotropic.
D-BI-18a consequently yields
\[
                       e(3-5B/a)\le12.                     \tag{UW.13}
\]
These arguments apply to the total degree and do not require
primitivity of the inducing block.

\emph{First, the large-parameter ranges.}
By (UW.2),
\[
 A-3B=\frac{(q-3)q^{n-1}-7q+3}{q+1}.                       \tag{UW.14}
\]
For $q\ge4$, $n\ge4$, this is at least
$q^2-q-6+9/(q+1)$, whose value at four is $39/5>6$.
Its forward difference is $2q-9/((q+1)(q+2))>0$.
For $n=3$, $q\ge7$, (UW.14) equals
$q^2-4q-3+6/(q+1)$, with value $75/4>6$ at seven
and positive forward difference
$2q-3-6/((q+1)(q+2))$.
Since $e\ge a\ge A$, these strict bounds contradict (UW.12).
Only $q=2,3$ and $(n,q)=(3,4),(3,5)$ remain.

\emph{Second, the nonselfdual constituents.}
For $q=3$ the degree is $a=(3^n-(-1)^n)/4$ and
$B=(3^{n-1}+5)/4\le a/2$ for all $n\ge3$.
The direct stabilizer estimate (UW.3) and the tensor-cycle
bound give $\nu(g)\ge e/4$ for every nonscalar inertia.
The same holds at $(n,q)=(3,4)$, where $(a,B)=(13,23/5)$,
and at $(n,q)=(3,5)$, where $(a,B)=(21,17/3)$.
Nonidentity scalar inertia also satisfies the ordinary
moving-rank bound. Equation (UW.11) would give $e\le6$,
contrary to the respective factor degrees, which are at least seven.

For $q=2$ and $n\ge5$,
\[
 a=(2^n-(-1)^n)/3,\quad B=(2^{n-1}+3)/3,\quad
 \gamma=(a-B)/(2a)<1/4.
\]
All nonscalar normalizers on the tensor product have support
at least $\gamma e$. The exact budget implies
\[
 e(4a-6B)/a\le6,\qquad
        4a-6B=\frac{2^n-4(-1)^n-18}{3}.                    \tag{UW.15}
\]
For $n\ge6$ the rightmost expression is at least14, so
$e\ge a$ is impossible. At $n=5$, $(a,B)=(11,19/3)$
and $\gamma=7/33$; (UW.15) forces $e\le11$, hence $t=1$.
Every genuine inertia then moves at least $\lceil7/3\rceil=3$
dimensions, giving $D\ge36>28=2e+6$. The remaining
nonselfdual $q=2$ constituents have $(n,a)=(4,5)$.

\emph{Third, the selfdual constituents.}
For $q=3$, $n\ge4$, use $a\ge A$ in
\[
                    3a-5B\ge3^{n-1}-17/2>12,
\]
contradicting (UW.13). At $(n,q)=(3,4)$ the selfdual
constituent has $a=12$ and $3a-5B=13>12$.
At $(n,q)=(3,5)$ the selfdual degrees are20 and21;
already $3\cdot20-5\cdot17/3=95/3>12$.
For $q=2$ the sole selfdual constituent is $W_0$.
For $n\ge7$,
\[
 3a_0-5B=\frac{2^{n-1}+6(-1)^n-15}{3}
                   \ge\frac{2^{n-1}-21}{3}>12,
\]
again contradicting (UW.13).

The remaining small cases are covered in full by (UW.7).
For $\operatorname{PSU}_3(3)$, every tensor normalizer has
the quarter bound, so $e\le6$. This leaves only the single
degree-six possibility; its exact integer inequality is
$12\lceil6/4\rceil=24>18$, also impossible.
For the degree22 constituent of $\operatorname{PSU}_6(2)$,
the quarter bound contradicts $e\ge22$ immediately.
For the selfdual degree10 constituent of $\operatorname{PSU}_5(2)$,
the full-normalizer fifth bound and the two nontrivial
multiplier costs give
\[
                   2e+6\ge2(e/2)+10(e/5)=3e.
\]
Thus $e\le6$, contrary to $e\ge10$.

Only $\operatorname{PSU}_4(2)$ and its degrees5 and6 remain.
The complete small-core theorem D-SPW-04b, which allows
every compatible outer and every factor permutation,
excludes their tensor products and gives $t=1$.
Its final clause removes the degree-six case in precisely
the stated non-Eisenstein nonsplit scope. Its degree-four
case cannot appear here because it is not an ordinary
unitary Weil constituent of $\operatorname{SU}_4(2)$.
The remaining degree-five case is the already isolated
rank-five inducing gate. This proves every assertion.
\end{proof}

\paragraph{Boundary of the result.}
No conclusion about non-Weil unitary representations follows.
No assertion about arbitrary mixed normal-core types or a
nontransitive family of factor labels is made. The result
eliminates an infinite family of possible sixfold suppliers
and identifies its exact existing small endpoint; it produces
no cycle and proves no additional period-map dominance.
